\color{unidarkgreen}{\subsection{Реле тока блокировки РПН}}
\color{black}

\subsubsection{Общие сведения}

\begin{enumerate}[label=\arabic{section}.\arabic{subsection}.\arabic{subsubsection}.\arabic*,  align=left, labelsep=4pt, leftmargin=0pt, itemindent=57pt]

\item
Предусмотрено реле, блокирующее автоматику РПН, когда ток перегрузки трансформатора
превышает регламентируемое для обмотки значение. По умолчанию данное реле подключено к
трансформаторам тока стороны ВН трансформатора. Порог по току срабатывания задается параметром
\textbf{«Iбл РПН»}.

\begin{figure}[H]
\centering
\includegraphics[width=0.8\textwidth,height=0.8\textheight,keepaspectratio]{\fbpath/011.1181 РПН/img/RPN.pdf}
\captionof{figure}{Функциональный блок РТ блокировки РПН Т}\label{pic:RPN}
\end{figure}


\small
\begin{longtable}{|>{\centering\arraybackslash}m{5.3cm}|>{\centering\arraybackslash}m{4.2cm}|>{\centering\arraybackslash}m{1cm}|>{\centering\arraybackslash}m{3.3cm}|>{\centering\arraybackslash}m{1.8cm}|}
% Первая страница: полный заголовок
\caption{Перечень уставочных параметров функции блокировки РПН\hfill\vspace{-0.5\baselineskip}}\label{desc:tbl1}\\
\hline
\rowcolor{gray!30}
Наименование уставки & Диапазон & Шаг & Значение по умолчанию & Единица измерения \\
\hline
\endfirsthead

% Промежуточные страницы: «Продолжение таблицы»
\caption*{\hspace{3pt}\emph{Продолжение таблицы \ref{desc:tbl1}}\hfill\ }\vspace{-0.5\baselineskip} \\
\hline
\rowcolor{gray!30}
Наименование уставки & Диапазон & Шаг & Значение по умолчанию & Единица измерения \\
\hline
\endhead

%===m>ДТЗ|-
%===m
\end{longtable}
\normalsize	

\end{enumerate}